\chapter{Testing and validation}
\label{vol11:ch08}

\epigraph{The test that proves only what the design team believed cannot test what the design team did not believe; the substantive test is the one that risks disconfirmation.}{}

\chapmeta{Half-life: medium-life. Archetypes: uncertainty, ethics. Exercise target: 25. Project track: hybrid.}

\noindent
Testing is the engineering activity of producing evidence about the design's actual behaviour under specified conditions. Validation is the activity of confirming that the design serves the user need (Chapter 2's distinction). The two activities consume substantial design-process effort and generate the evidence that supports certification, market introduction, and ongoing service. This chapter develops the discipline: how tests are designed, what test types serve which purposes, how statistical and accelerated methods extend test reach, how software testing differs from hardware testing, how field tests integrate with structured testing, how the data are managed, and how test plans fail when they prove only what the design team already believed.

\section{Test-design as engineering discipline}\pathtag{core}

A test plan is itself a design problem. The plan specifies what will be tested, under what conditions, with what measurements, against what acceptance criteria, with what statistical analysis. The plan is reviewed before testing begins; the testing executes the plan; the results are analysed against the plan; the plan's adequacy is itself reviewed at the end.

The elements of a substantive test plan:

\textbf{Test objectives.} What the test is intended to learn or to confirm. Objectives must be specific (the test produces measurable findings against them) and complete (the objectives cover the design's verification needs). "Test the design" is not an objective; "verify that the design meets requirements R1.1, R1.2, and R2.3 under operating conditions $-20$ to $+50\,\degC$" is an objective.

\textbf{Test items.} The specific designs (or prototypes; or production units) that will be tested. The test items must be representative of the design population the conclusions will be applied to; testing a hand-built prototype and inferring production behaviour is non-conservative.

\textbf{Test conditions.} The conditions under which the items are tested: inputs (electrical, mechanical, thermal, environmental); load profiles; duration; sequence. The conditions must be controlled (the test produces interpretable results) and representative (the conclusions apply to service conditions).

\textbf{Measurements.} The data captured during the test: what is measured; with what instruments; with what precision and accuracy; with what sampling rate; with what storage and traceability. The measurements determine what the test can show; data that are not collected cannot be analysed.

\textbf{Acceptance criteria.} The thresholds at which the test passes or fails. The criteria must be quantitative (the pass/fail decision is decidable) and linked to the requirements (the criteria express the requirement's intent). "The design must perform well" is not a criterion; "the output must remain within $\pm 5\%$ of nominal across the operating range" is a criterion.

\textbf{Statistical plan.} For tests that sample from a population (e.g., reliability testing of $N$ units), the statistical analysis is part of the plan: the sample size needed to achieve the desired confidence; the analysis method; the criteria for accepting the population claim. The statistical plan prevents post-hoc analysis from inflating claims beyond what the test data support.

\textbf{Test schedule.} The sequence and timing of tests; the dependencies between tests; the resources required. The schedule supports the project's design-process schedule; tests scheduled too late produce findings that are too late to act on.

\textbf{Roles and responsibilities.} Who designs the test, who executes it, who analyses the data, who reviews the findings. The separation of roles (the test executor is not the design's author; the analyst is not the test executor) reduces the confirmation-bias risk.

\textbf{Documentation requirements.} What the test produces in writing: the test plan; the test procedures; the test report; the data records. The documentation is the persistent form of the test's findings.

The discipline of test-plan design is the discipline of producing plans that are specific, testable, and that risk disconfirmation. Plans that are vague (so any outcome can be interpreted as success), or that are not risky (so the design is certain to pass), are theatre. Substantive test plans surface design issues that the design team did not anticipate; the design team's response to those issues is the substantive engineering work.

\begin{principle}[Tests that risk disconfirmation]
A test that the design is certain to pass is not a substantive test; it is a demonstration. The substantive test is the test that risks disconfirmation: where the design might fail, and where the failure would identify an issue the design team needs to address. The discipline of test design is the discipline of producing plans whose execution can fail informatively.
\end{principle}

\section{Functional, performance, environmental, life testing}\pathtag{core}

Test types serve different purposes; the test programme typically includes multiple types:

\textbf{Functional testing.} The design's intended functions are exercised; each function is tested for correct behaviour. The testing covers the specified functions individually and in combination. For a software product: each feature is tested; the integration of features is tested. For a mechanical assembly: each mechanism is tested; the interactions are tested.

\textbf{Performance testing.} The design's quantitative performance is measured against the requirements. For a vehicle: acceleration; braking distance; fuel consumption; emissions. For a server: throughput; response time; concurrency; latency at percentiles. Performance testing produces the numbers that go into the product's specification documents.

\textbf{Environmental testing.} The design is exercised under the environmental conditions of its intended service: temperature extremes (high and low); humidity; vibration; shock; thermal cycling; salt fog (for marine applications); altitude (for aerospace); EMI/EMC (for electronics). The testing identifies environmental sensitivities and verifies the design meets its environmental requirements.

\textbf{Life testing.} The design is operated for extended periods (sometimes accelerated; see Section 4) to characterise its wear-out behaviour, failure modes, and service life. Life testing produces the data that support reliability claims and maintenance-interval specifications.

\textbf{Reliability testing.} The design's failure rate over a specified operating period is measured. Reliability testing typically uses multiple units (the population for which the reliability claim is made) and may use accelerated stress to compress the testing duration.

\textbf{Stress testing.} The design is exercised at conditions exceeding its specification to identify failure modes and margins. Stress testing produces the evidence that the design has the safety margins specified; it also surfaces failure modes the at-specification testing does not.

\textbf{Compatibility testing.} The design is tested with the systems and components it must interface with: electrical compatibility with power sources; mechanical compatibility with mating parts; protocol compatibility with peer devices. Compatibility testing is particularly important for designs that integrate into larger systems (consumer products with appliances; medical devices with hospital infrastructure).

\textbf{Usability testing.} The design is operated by representative users; their behaviour, errors, and feedback are documented. Usability testing produces the evidence about the design's user-facing behaviour that engineering analysis alone cannot.

\textbf{Safety testing.} The design's behaviour under fault conditions, abuse conditions, and emergency conditions is verified. Safety testing demonstrates that the design fails safely, that the protections operate as designed, and that the design satisfies regulatory safety requirements.

\textbf{Regulatory and certification testing.} The design is tested against the specific procedures and conditions the regulatory framework requires. The testing may be conducted by the design organisation, by independent test laboratories, or by the regulator depending on the regulatory regime. The results feed the certification submission.

The discipline of test programme design is the discipline of selecting the types and their depths to match the design's risks. A safety-critical aerospace component requires extensive testing across all types; a low-stakes consumer product may require functional and basic environmental testing only. The mismatch (testing a low-stakes product to safety-critical depth wastes resources; testing a safety-critical product to consumer-product depth produces inadequate evidence) is the substantive engineering judgment.

\section{Statistical design of experiments}\pathtag{standard}

Many tests are statistical: the design's behaviour varies across instances; the test samples from the variation; the conclusion is a statistical claim about the population. The discipline of statistical design of experiments (DOE) addresses the test-plan design problem for these tests.

The classical DOE techniques:

\textbf{Factorial design.} The test varies multiple factors simultaneously; the response is measured at each factor combination. A full factorial with $k$ factors at $n$ levels requires $n^k$ runs; the technique characterises the main effects and the interactions but the cost scales steeply with $k$.

\textbf{Fractional factorial.} A subset of the full factorial is run; the subset is designed to characterise the main effects with limited interaction characterisation. The technique reduces the run count substantially at the cost of confounding some interactions; the confounding is documented and the design is chosen to confound only interactions that are expected to be small.

\textbf{Response-surface methodology.} The factor space is sampled at a structured pattern (central composite design; Box-Behnken design); the response is modelled as a polynomial; the model is used to predict the response across the factor space and to identify optimal factor settings.

\textbf{Taguchi methods.} The factors are categorised as control factors (the engineer sets them) and noise factors (the engineer cannot directly control them in service); the experimental design tests the response across both, identifying factor settings that minimise sensitivity to noise. The technique is dominant in manufacturing-process robustness.

\textbf{Latin hypercube sampling.} The factor space is sampled stratified across each factor; the sampling supports surrogate-model construction for high-dimensional factor spaces.

\textbf{Sequential designs.} The experiment is conducted in stages; the early-stage results inform the design of the later stages. The technique uses the experimental budget efficiently by concentrating effort on the regions of factor space that the early results identified as informative.

The DOE framework also addresses:

\textbf{Sample size determination.} How many test units are needed to detect an effect of specified size with specified power? The sample-size formulas depend on the statistical model and the desired confidence; underpowered tests cannot distinguish a real effect from noise; overpowered tests waste resources.

\textbf{Randomisation.} The order of test runs is randomised to prevent systematic biases from time-varying factors (temperature drift; operator fatigue; equipment aging). Non-randomised tests can produce spurious correlations between the run order and the response.

\textbf{Replication.} Each combination of factors is tested multiple times to characterise the within-condition variability. Replication separates the variability due to factor changes from the variability inherent to the test conditions.

\textbf{Blocking.} Known sources of variation that are not the experiment's focus are blocked: e.g., units tested on the same day are blocked together. Blocking removes the blocked variation from the residual error, improving the experiment's power.

The discipline of DOE is taught in statistical-engineering programs and is codified in standards (ASTM E1325 for terminology; ASTM E1488 for statistical procedures). The major statistical software packages (Minitab; JMP; R; Python's scipy and statsmodels) support DOE design and analysis.

\section{Accelerated testing}\pathtag{standard}

Many tests address behaviour over service lives that exceed the available test time. A design with a 20-year service life cannot be tested at full duration before commitment; the test must be accelerated to produce evidence in available time.

Accelerated testing relies on a model relating the test conditions to the service conditions: the stress at test exceeds the stress in service; the time to failure scales by an acceleration factor that depends on the stress model.

The common models:

\textbf{Arrhenius for thermal acceleration.} Failure rate scales with $\exp(-E_a / kT)$ where $E_a$ is the activation energy and $T$ is absolute temperature. Increasing test temperature accelerates the failure mechanisms whose rate is dominated by Arrhenius behaviour (many semiconductor failure mechanisms; some chemical degradation processes). The acceleration factor is computed from the activation energy and the test-to-service temperature difference.

\textbf{Voltage acceleration.} For electrical failure mechanisms, failure rate scales with $V^n$ where $n$ is empirically determined. The voltage acceleration applies to dielectric breakdown, electrochemical migration, and similar mechanisms.

\textbf{Coffin-Manson for thermal cycling.} For thermomechanical fatigue from temperature cycling, the number of cycles to failure scales with $(\Delta T)^{-n}$ where $\Delta T$ is the cycle's temperature swing. The exponent $n$ depends on the failure mechanism and the material.

\textbf{S-N for mechanical fatigue.} For mechanical fatigue, the number of cycles to failure scales with the stress amplitude through the material's S-N curve (Volume X Chapter 2). Increased load amplitude accelerates the fatigue testing.

\textbf{Power-law models for general acceleration.} For mechanisms that do not fit a specific physical model, the acceleration factor is sometimes modelled as a power law in a stress quantity; the exponent is determined experimentally.

\textbf{HALT/HASS (Highly Accelerated Life Test / Highly Accelerated Stress Screen).} Test protocols that combine extreme thermal cycling and vibration to surface failure modes rapidly. HALT is used in development to identify weakness; HASS is used in production to screen units that have specific defects.

The discipline of accelerated testing includes:

\textbf{Acceleration-factor validation.} The acceleration factor depends on the mechanism; testing must include mechanism identification (e.g., post-test failure analysis to confirm that the failure modes at test are the same as those expected in service). Acceleration factors that are correct for one mechanism are wrong for another; testing that accelerates a non-service-relevant mechanism produces misleading results.

\textbf{Stress-limit awareness.} Acceleration can be pushed only to the point where the accelerated conditions still produce the same failure mechanisms as the service conditions. Beyond that limit, the test surfaces different mechanisms that do not predict service behaviour. The discipline includes characterising the stress-mechanism relationship and respecting the limit.

\textbf{Statistical analysis at small sample.} Accelerated tests typically use small samples (high per-unit cost; high per-test cost); the statistical analysis must account for the small sample (Weibull analysis; lognormal analysis; the various small-sample reliability techniques).

\textbf{Validation against service.} Where possible, the accelerated-test predictions are validated against service data from prior generations of the design. Validation evidence supports the application of the acceleration model to new designs.

\section{Software testing revisited from Volume VII}\pathtag{standard}

Software testing differs structurally from hardware testing: the software's behaviour is determined by its code rather than by its physical properties; the code can be tested exhaustively for small input spaces but not for large ones; the testing can be automated to a degree hardware testing cannot. Volume VII developed software testing as part of the information-engineering chapters; this section revisits it as part of the design's testing programme.

The dominant techniques:

\textbf{Unit testing.} Each function or method is tested in isolation; the tests are typically written by the developer and run automatically. Unit tests catch implementation errors; their coverage (the fraction of the code's branches that the tests exercise) is a quality metric.

\textbf{Integration testing.} The interaction of multiple units is tested; integration tests catch interface mismatches and emergent behaviours.

\textbf{System testing.} The complete system is tested against the requirements; system tests verify the integrated software meets its functional and non-functional requirements.

\textbf{Acceptance testing.} The user (or the user's representative) tests the system against their needs; acceptance tests confirm the system serves the user need.

\textbf{Regression testing.} The full test suite is rerun after changes; regression tests catch changes that broke previously-working behaviour. Automated regression suites are the foundation of continuous-integration / continuous-deployment software development.

\textbf{Property-based testing.} The test specifies properties the software should satisfy; the framework generates many random inputs and checks the properties. The technique surfaces edge cases that example-based testing would miss.

\textbf{Fuzzing.} The software is exercised with random or semi-structured inputs to discover crashes and security vulnerabilities. Modern fuzzers (AFL; libfuzzer; coverage-guided fuzzers) are highly effective at finding bugs in parsers, decoders, and security-relevant code.

\textbf{Formal verification.} The software's properties are proven mathematically; the proof is the evidence of correctness. The technique is expensive but produces the strongest evidence available; it is used in safety-critical and security-critical domains where the cost is justified.

\textbf{Static analysis.} The software is analysed without execution; the analysis identifies potential bugs (null dereferences; buffer overflows; race conditions; type errors). Modern static analysers (Coverity; Klocwork; Pylint; the various sanitisers and compilers' warning systems) catch many bugs early.

\textbf{Performance testing.} The software's performance under load is measured; the testing identifies bottlenecks, scaling limits, and resource consumption.

The discipline of software testing includes:

\textbf{Test-driven development.} The test is written before the implementation; the test specifies the expected behaviour; the implementation is developed to pass the test. The technique produces high test coverage and resists implementation drift from the specification.

\textbf{Coverage analysis.} Statement coverage, branch coverage, path coverage, and modified condition/decision coverage (MC/DC, required by avionics standards) measure the testing's reach. Coverage is a necessary but not sufficient condition for adequate testing; high coverage with weak assertions does not catch bugs the assertions would have caught.

\textbf{Mutation testing.} Small changes (mutations) are introduced into the code; the test suite should detect each mutation. Mutations that are not detected indicate gaps in the test suite. The technique evaluates the test suite's substantive strength.

\textbf{Continuous integration.} Every code change is tested automatically; the test results are reported to the developers immediately; failed tests block the change from integration. The technique catches integration issues at the moment they are introduced rather than late in the development cycle.

The Therac-25 case (Volume X Chapter 5) is partly a software-testing failure: the race condition in the operator-interface code was not surfaced by the testing because the testing did not exercise the fast-typing input sequence that triggered the race. The software testing's coverage was inadequate to the failure mode; the system testing did not include the conditions that the fast-typing operators routinely encountered.

\section{Field testing and beta programmes}\pathtag{standard}

Field testing exposes the design to actual operating conditions in actual user hands. The testing surfaces issues that controlled testing cannot:

\textbf{Population diversity.} The user population in the field is more diverse than any test population; behaviours that the test population did not exhibit may be common in the field.

\textbf{Environmental diversity.} The field environments cover the full operating envelope; controlled testing covers a sample. Environmental conditions that the controlled testing did not include may produce field issues.

\textbf{Use-case diversity.} Field use surfaces use cases the design team did not anticipate; some are creative (the user repurposes the design for an unanticipated need that the design serves well), and some are problematic (the user uses the design in ways the design does not handle).

\textbf{Long-duration behaviour.} The field testing accumulates operating hours faster than controlled testing; the accumulation surfaces long-duration behaviours that compressed testing did not.

\textbf{Integration with user workflows.} The design's integration with the user's existing tools, processes, and habits is tested only in the field; controlled testing typically isolates the design from the user's full context.

The structures for field testing:

\textbf{Limited customer release.} A small set of customers receives the product before broad release; their use produces field data and feedback.

\textbf{Beta programmes.} Selected users receive pre-production units; the relationship is structured (beta-test agreements; defined feedback channels); the users accept the early-version risks in exchange for early access.

\textbf{Canary releases.} For software, a new version is released to a small fraction of users first; the deployment is expanded if no issues surface, rolled back if issues do.

\textbf{Phased rollouts.} The release is geographically or demographically staged; the staging allows issues to be detected at small scale before broad exposure.

\textbf{Telemetry and analytics.} The deployed units send back data about their use, performance, and failures; the data feed design improvement and operational decisions.

The field testing's value depends on the data flowing back to the design team. Beta programmes that ship units without structured feedback collection produce little engineering value; programmes with structured feedback produce substantial input to subsequent design iterations.

\section{Test-data management and reproducibility}\pathtag{standard}

The test data are an engineering asset. The data support the design's certification, its forensic analysis after service failures, the next-generation design's development, the comparison with the design's actual service behaviour. The data's value depends on their preservation, accessibility, and traceability.

The discipline of test-data management:

\textbf{Data collection standards.} The data are collected at specified rates, with specified precision, in specified formats. The standards are documented; deviation from the standards is documented; the data's quality is assessed against the standards.

\textbf{Provenance.} Each data record is linked to its source: the test article (serial number); the test conditions (date, environment, instrument calibration); the operator; the data acquisition system version. Provenance supports the data's interpretation and the audit trail.

\textbf{Storage and access.} The data are stored in systems that support long-term retention (decades for safety-critical applications; multi-decades for nuclear and aerospace); access is controlled; access logs preserve who accessed what.

\textbf{Reproducibility.} The test conditions are documented to the level required for the test to be reproducible by an independent analyst. The data analysis is documented likewise; the analysis can be reproduced from the raw data.

\textbf{Version control.} The data, the analyses, the test procedures, and the test reports are version-controlled. Changes are tracked; the relationships between versions are explicit.

\textbf{Quality flags.} The data records carry quality flags (raw; preprocessed; validated; suspect; rejected) that subsequent analysis uses to weight or filter the data. Quality flags preserve the data's nuance without losing the records.

\textbf{Integration with the design record.} The test data are linked to the design elements they test; the linkage supports requirements traceability (Chapter 2) and certification (Chapter 9).

The cost of test-data management is real (storage; processing; documentation; access infrastructure); the cost of inadequate management is the cost of test data that cannot be found, cannot be interpreted, or cannot be trusted when they are needed. The discipline of test-data management is part of the engineering organisation's infrastructure; it is built over years and is destroyed quickly under cost pressure that does not recognise the data's value.

\begin{archetype}[uncertainty]
Testing is the uncertainty archetype's engineering response. The design's actual behaviour is uncertain; testing produces evidence that reduces the uncertainty. The discipline of testing is the discipline of selecting the tests that produce the most uncertainty reduction per resource, of conducting the tests with the discipline that makes the evidence trustworthy, and of preserving the evidence so subsequent decisions can use it.
\end{archetype}

\section{Failure: test plans that proved only what they were designed to prove}\pathtag{core}

The failure mode this section addresses is the test plan that confirms the design's intended behaviour but does not test the behaviour the design might exhibit. The test runs; the design passes; the design is committed; in service, the design exhibits behaviours the test did not address; the failures occur.

The mechanisms:

\textbf{Limited operating envelope.} The test covers the operating conditions the design team expected; the design encounters conditions in service that the test did not include. The 737 MAX MCAS testing covered MCAS activation under conditions where MCAS was working correctly; it did not adequately test MCAS activation under sensor failure conditions, which is the condition under which the accidents occurred \cite{acc:house-737max-2020}.

\textbf{Single-user-class testing.} The usability testing covered users similar to the design team; the design encounters users in service whose interaction patterns differ; the design fails for the unaccounted users. The Therac-25 testing covered operators with the original-machine experience; the field operators with the new machine's faster workflow encountered the race condition the testing did not surface \cite{paper:leveson-turner1993}.

\textbf{Component-level only.} The testing covers each component individually; the integrated system's behaviour is not tested; integration-level failures occur in service. The Ariane 5 Flight 501 failure (Volume X Chapter 4) had elements: the software components had been tested individually but the integration with the Ariane 5's flight envelope (which differed from Ariane 4's) had not been adequately tested \cite{acc:ariane501-ifr}.

\textbf{Steady-state only.} The testing covers steady-state behaviour; transient and edge-condition behaviours are not tested; failures occur during transients. Many control-system failures have this character: the steady-state response is well-characterised; the transient response under unusual conditions is not.

\textbf{Pass/fail without margin.} The test confirms the design meets specifications without characterising the margin to failure. A design that just barely passes is treated as equivalent to a design that comfortably passes; the margin difference matters in service.

\textbf{Confirmation-only assertions.} The test assertions check that the design produces the expected output; assertions about unexpected outputs (the design produces output X but should have produced output Y) are not present. The asymmetric assertions confirm what the design team believed; they do not catch what the design team did not anticipate.

\textbf{No negative tests.} The test plan includes tests for the design's correct response to correct inputs; it does not include tests for the design's response to incorrect inputs (out-of-range; malformed; adversarial). The design's failure mode under negative inputs is not characterised.

\textbf{Test conducted by the design team.} The test was designed and executed by the people who designed the system; confirmation bias shapes the test toward outcomes the design team expects. Independent test design and execution resists this bias.

\textbf{Acceptance-criteria gaming.} The acceptance criteria are set after preliminary testing; the criteria are calibrated to ensure the design passes. The substantive engineering question (does the design's performance meet what the user needs?) is replaced by the procedural question (does the design pass the criteria the design organisation set?).

\textbf{Statistical inadequacy.} The sample size is too small to detect the effect; the test passes by failing to detect what it should have detected. The Volkswagen emissions case has elements: the test conditions did not exercise the engine under conditions the engine's defeat-device logic responded to differently \cite{acc:epa-vw-nov-2015}; the test result was correct for the test conditions, wrong for the service conditions.

The remediation is the discipline of test-plan design that risks disconfirmation (Section 1), independent test design and execution, statistically-adequate samples, negative testing alongside positive, integrated-system testing alongside component testing, and operating-envelope coverage that includes the conditions the design might encounter rather than only the conditions the design team expected.

\begin{principle}[Test what the design might do, not only what it should do]
The substantive test programme includes tests for what the design might do (under conditions the design team did not fully anticipate) alongside tests for what the design should do (under conditions the design team expected). Test plans that only confirm the expected behaviour produce evidence that fails to surface the design's actual failure modes; the failures occur in service rather than in test.
\end{principle}

\begin{project}[Test plan and a manufacturer-missed condition]
Select one product the reader uses regularly (a consumer device; a software tool; an appliance). Design a test plan for the product's principal function. The plan should cover: functional verification; performance characterisation; environmental envelope; failure modes; usability under realistic conditions. Then identify one specific condition or behaviour that the manufacturer's actual testing appears to have missed (evidenced by the product's failures, user complaints, or known limitations). Identify the test that would have caught it and the reason it was not in the manufacturer's test plan. Deliverable: a 4-page document with the test plan and the missed-condition analysis.

\textbf{Hazard.} The exercise often reveals that the product the reader uses has substantial test gaps that the user has been quietly absorbing. The point is the engineering insight, not the consumer complaint; the document's value is in identifying what the test discipline would catch versus what it has missed.
\end{project}

\section{Exercises}

\subsection*{Calculation}\pathtag{core}

\begin{exercise}
A reliability test must demonstrate that the design's failure rate is below $10^{-4}$ per hour with 90\% confidence. Using the exponential distribution and the chi-squared lower bound formula, compute the test-hours required for: $n = 1$ failure during test; $n = 0$ failures (success); $n = 5$ failures. Discuss the trade-off between test duration and acceptable test failures.
\end{exercise}

\begin{exercise}
A full factorial test of 5 factors at 3 levels requires $3^5 = 243$ runs; each run costs \$2{,}000. Compute the total cost. A fractional factorial (resolution V; $3^{5-1} = 81$ runs) costs how much; what main effects and two-factor interactions can it estimate?
\end{exercise}

\begin{exercise}
An Arrhenius acceleration factor is computed for a service temperature of $40\,\degC$ and a test temperature of $85\,\degC$, with activation energy $0.7\,\text{eV}$. Compute the acceleration factor. If service life is 10 years, compute the equivalent test duration.
\end{exercise}

\begin{exercise}
A software test suite achieves 80\% branch coverage. The remaining 20\% includes the most rarely-executed code paths. Compute the expected number of defects in the uncovered code, given a defect density of $1$ defect per $1{,}000$ lines and $50{,}000$ lines of code with the rare paths comprising 20\% of total lines.
\end{exercise}

\begin{exercise}
A canary release deploys to 1\% of users for 24 hours before broad rollout. The deployment population is 10 million users. Compute the canary population. If the design has a defect that affects 0.1\% of users, compute the expected number of canary users who would experience the defect and discuss whether 24 hours is sufficient for detection.
\end{exercise}

\subsection*{Derivation}\pathtag{standard}

\begin{exercise}
Derive the sample size required to detect an effect of size $\Delta$ in a means-comparison test with significance $\alpha$ and power $1-\beta$. Identify the sample-size scaling with $\Delta$ and discuss the practical implications for testing small effects.
\end{exercise}

\begin{exercise}
For an accelerated-life test with multiple stress levels, derive the maximum-likelihood estimator for the activation energy and the median life at service conditions. Identify the conditions under which the estimator is unbiased.
\end{exercise}

\begin{exercise}
Derive the confidence interval for the failure rate when zero failures are observed in a test of duration $T$. Identify the difference between the classical (frequentist) and Bayesian (with uninformative prior) intervals.
\end{exercise}

\subsection*{Estimation}\pathtag{standard}

\begin{exercise}
Estimate the test-engineering hours required for the certification of a commercial aircraft model. Compare to the design-engineering hours. Discuss the implication for project staffing.
\end{exercise}

\begin{exercise}
Estimate the cost of a complete environmental-test programme for a consumer electronic product (temperature, humidity, vibration, shock, EMI, ESD, IP rating). Identify the dominant cost categories.
\end{exercise}

\begin{exercise}
Estimate the fraction of major product recalls that are attributable to inadequate testing versus inadequate design versus inadequate manufacturing. Discuss what the fractions imply about engineering investment priorities.
\end{exercise}

\begin{exercise}
Estimate the time required to develop a comprehensive automated test suite for a typical software product (e.g., 100{,}000 lines of code, web application). Identify the relationship between test development time and code development time.
\end{exercise}

\subsection*{Simulation}\pathtag{standard}

\begin{exercise}
Implement a Monte Carlo simulation of a fractional factorial design. Generate synthetic data with known main effects and interactions; analyse with the fractional design; identify the conditions under which the design correctly identifies the main effects and the conditions under which confounding obscures them.
\end{exercise}

\begin{exercise}
Simulate an accelerated-life test: units fail according to a Weibull distribution at the service stress; the test is conducted at higher stress with the failures scaled by an acceleration factor; the analyst estimates the service-condition life from the test data. Sweep the test stress and identify the conditions under which the estimate is accurate versus biased.
\end{exercise}

\begin{exercise}
Implement a test-coverage simulation: a code base has 1000 branches; the test suite has $N$ test cases; each test case exercises a random subset of branches. Sweep $N$ and identify the relationship between test count and coverage; identify the saturation point and the cost of approaching 100\% coverage.
\end{exercise}

\subsection*{Design}\pathtag{standard}

\begin{exercise}
Design a test plan for a new medical infusion pump (Chapter 3 example). Cover: functional testing; performance testing; environmental testing; safety testing; usability testing; reliability testing; regulatory testing. Identify the sequencing, the resources required, and the integration with the design schedule.
\end{exercise}

\begin{exercise}
Design a statistical design of experiments for optimising a 3D-printing process (factors: layer height, print speed, temperature, infill density; response: part strength). Specify the design type, the runs, the analysis approach, and the optimisation.
\end{exercise}

\begin{exercise}
Design an accelerated-life test programme for an automotive battery pack. Specify the acceleration models, the stress conditions, the duration, the sample size, and the validation against service data.
\end{exercise}

\begin{exercise}
Design a beta programme for a new software product. Specify the beta-user selection, the feedback mechanisms, the data telemetry, the success criteria for general release, and the procedures for handling beta-discovered issues.
\end{exercise}

\subsection*{Judgement}\pathtag{mastery}

\begin{exercise}
Argue: the 737 MAX MCAS certification testing failed because it tested what the design team believed MCAS would do rather than what MCAS might do under sensor failure. Discuss the certification-testing changes that would have addressed the failure mode; identify the institutional structures that would compel such changes.
\end{exercise}

\begin{exercise}
Discuss the Volkswagen diesel emissions case as a test-plan failure. The certification test was conducted as specified; the test result was technically correct; the design's actual service behaviour diverged. Identify the testing-discipline changes that would have caught the divergence and the regulatory-framework changes that have been made in response.
\end{exercise}

\begin{exercise}
Argue: the engineering profession's reliance on accelerated testing for service-life prediction is structurally fragile; service-life predictions from accelerated testing have known failure modes that are repeatedly missed. Identify the conditions under which accelerated testing is reliable versus unreliable.
\end{exercise}

\begin{exercise}
A test engineer is asked by management to revise the test plan so the design passes. Discuss the engineer's options under the professional-ethics frameworks (NSPE; IEEE; ASME). Identify the institutional structures that protect the test engineer from the management pressure and the conditions under which the structures fail.
\end{exercise}

\begin{exercise}
Discuss the structural reasons that field testing and beta programmes often produce less engineering value than they could. Identify the institutional commitments required to make the programmes substantive rather than theatre.
\end{exercise}

\begin{exercise}
A test produces a result that contradicts the design team's expectation. Three responses are possible: (a) revise the design; (b) revise the test (it must have been wrong); (c) revise the expectation (the design behaves differently than assumed). Discuss the conditions under which each response is appropriate and the institutional structures that protect against systematically choosing the response that does not require design revision.
\end{exercise}
